\section*{PHỤ LỤC}
\phantomsection\addcontentsline{toc}{section}{\numberline{} PHỤ LỤC}

\subsection*{A. Tài nguyên dự án}

\begin{itemize}
    \item \textbf{Mã nguồn (GitHub):}
    \url{https://github.com/NLT22/multi\_stream\_people\_tracker}

    \item \textbf{Video demo và hình ảnh (Google Drive):}
    \url{https://drive.google.com/drive/folders/1yO3T-kCAQ-KZG0K4Tydum4X66aGpdzhx?usp=sharing}
\end{itemize}

\subsection*{B. Cấu trúc thư mục demo}

Thư mục demo được tổ chức theo từng môi trường. Mỗi môi trường có một bộ đầy đủ gồm 3 video và thư mục heatmap:

\begin{center}
\texttt{output/demo/<môi-trường>/}
\end{center}

\begin{table}[H]
\centering
\begin{tabular}{|l|p{10cm}|}
\hline
\bfseries File & \bfseries Nội dung \\
\hline
\texttt{<env>\_live\_buffered\_osd.mp4} & 4 camera đồng thời, hiển thị Buffered ID ổn định + quỹ đạo di chuyển \\
\texttt{<env>\_tracking\_bev.mp4}       & Góc nhìn từ trên xuống (BEV): mỗi người một chấm màu + vệt đường đi \\
\texttt{<env>\_analytics.mp4}           & ROI đếm người + cảnh báo quá tải + vạch đếm vào/ra (cấu hình riêng từng môi trường) \\
\texttt{heatmap/}                       & 15 ảnh PNG: \texttt{cam\_\{0--3\}\_\{typeheatmap\}.png} và \texttt{bev\_\{typeheatmap\}.png} \\
\hline
\end{tabular}
\end{table}

Năm môi trường: \texttt{office}, \texttt{cafe\_shop}, \texttt{lobby}, \texttt{industry\_safety}, \texttt{retail}. Tất cả video sử dụng chế độ \textbf{Buffered ID (anchor-guided)} để hiển thị định danh ổn định.
